% !TEX program = xelatex
% 编译方式:在 Overleaf 中将编译器设置为 XeLaTeX
\documentclass[11pt,a4paper]{article}

% ============ 宏包 ============
\usepackage[UTF8]{ctex}               % 中文支持
\usepackage[left=1.8cm,right=1.8cm,top=1.8cm,bottom=1.8cm]{geometry}
\usepackage{xcolor}
\usepackage{graphicx}
\usepackage{enumitem}
\usepackage{titlesec}
\usepackage{tikz}
\usepackage{fontawesome5}             % 图标(可选)
\usepackage{hyperref}
\usepackage{array}
\usepackage{tabularx}

% ============ 颜色定义 ============
\definecolor{themecolor}{RGB}{196, 30, 58}   % 红色主题
\definecolor{tagcolor}{RGB}{196, 30, 58}
\definecolor{textgray}{RGB}{80, 80, 80}

% ============ 超链接样式 ============
\hypersetup{
    colorlinks=true,
    urlcolor=themecolor,
    linkcolor=themecolor
}

% ============ 段落设置 ============
\setlength{\parindent}{0pt}
\pagestyle{empty}

% ============ 章节标题样式 ============
\titleformat{\section}
    {\color{white}\bfseries\large}
    {}
    {0pt}
    {\colorbox{themecolor}{\parbox{\dimexpr\textwidth-2\fboxsep}{\strut #1}}}
\titlespacing*{\section}{0pt}{12pt}{6pt}

% ============ 小标签(C9、985) ============
\newcommand{\tag}[1]{%
    \tikz[baseline=(X.base)]\node[rectangle, rounded corners=2pt,
    fill=tagcolor, text=white, inner sep=2pt, font=\footnotesize\bfseries] (X) {#1};%
}

% ============ 条目标题(左标题 + 右日期) ============
\newcommand{\entry}[2]{%
    \noindent\textbf{#1}\hfill{\color{textgray}#2}\\
}

\newcommand{\subentry}[2]{%
    \noindent#1\hfill{\color{textgray}#2}\\
}

% ============ 列表样式 ============
\setlist[itemize]{leftmargin=1.2em, itemsep=2pt, topsep=2pt, parsep=0pt}

% ============ 开始文档 ============
\begin{document}

% ====================== 个人信息 ======================
\begin{minipage}[t]{0.75\textwidth}
    \vspace{0pt}
    {\Huge\bfseries\color{themecolor} 李金美}\\[8pt]
    \faPhone\ 18790566306 \quad \faEnvelope\ 18790566306@163.com\\[4pt]
    23岁 \ | \ 女 \ | \ 籍贯:河南新乡 \ | \ 汉族 \ | \ 中共党员
\end{minipage}%
\hfill
\begin{minipage}[t]{0.22\textwidth}
    \vspace{0pt}
    \begin{flushright}
    \framebox[3cm]{\rule{0pt}{3.6cm}\hspace{-3cm}\parbox{3cm}{\centering \small 照片\\(可替换)}}
    \end{flushright}
\end{minipage}

\vspace{6pt}

% ====================== 教育背景 ======================
\section{教育背景}

\entry{中国科学技术大学 \tag{C9} \tag{985}}{2025年09月 - 至今}
\subentry{网络空间安全专业 \ 硕士 \ 网络空间安全学院}{合肥}
\begin{itemize}
    \item \textbf{保研学硕,中国科学技术大学硕士一等学业奖学金}
\end{itemize}

\vspace{4pt}

\entry{武汉大学 \tag{985}}{2021年09月 - 2025年06月}
\subentry{网络空间安全专业 \ 本科 \ 国家网络安全学院}{武汉}
\begin{itemize}
    \item \textbf{GPA:3.92/4.0(前3\%)}
    \item \textbf{主修课程:}编译原理、计算机网络、计算机组成原理、数据结构、操作系统原理及安全、密码学、数据库系统及安全等
    \item \textbf{奖项荣誉:}国家奖学金、武汉大学优秀本科毕业生、雷军计算机创新与发展基金、武汉大学三好学生、武汉大学优秀学生奖学金
    \item \textbf{科研成果:}发明专利《基于声学传感的智能用户认证方法、设备及软件产品》第一作者,授权专利号:ZL 202410993112.5;2024年第十七届全国大学生信息安全竞赛(A类学科竞赛)国家级一等奖;2024年第十七届中国大学生计算机设计大赛(A类学科竞赛)中南地区赛一等奖
\end{itemize}

% ====================== 实习经历 ======================
\section{实习经历}

\entry{科大讯飞股份有限公司 ---- 讯飞医疗大模型实习}{2025年12月 - 2026年06月}

% -------- 环境搭建 --------
\textbf{H200 多机 Post-training 训练环境从零搭建（SFT + GRPO）}
\begin{itemize}
    \item 从零完成 H200 多机集群 + qwen3.5 post-training 训练环境部署和调试（verl / Megatron-LM / vLLM / Ray / Swift）；
    \item 调试解决训练过程中遇到的报错bug，如训练随机崩溃等，确保训练环境和脚本稳定可复用。
\end{itemize}

\textbf{{SFT 医学数据全量微调与分析（Qwen3.5-27B）}}
\begin{itemize}
    \item 基于 Swift 框架（Megatron后端）完成 Qwen3.5-27B 全量 SFT 训练（230 万条医学数据）；对 SFT 后 IFEval 指令遵循能力下降进行归因分析，定位为思维链缩短与训练数据分布偏移，提供改进建议。
\end{itemize}
\textbf{{SFT 医学数据全量微调与分析（Qwen3.5-27B）}}
\begin{itemize}
    \item 基于 Swift 框架（Megatron后端）完成 Qwen3.5-27B 全量 SFT 训练（230 万条医学数据）；对 SFT 后 IFEval 指令遵循能力下降进行归因分析，定位为思维链缩短与训练数据分布偏移，提供改进建议。
\end{itemize}

% -------- 病历质控 GRPO --------
\textbf{病历质控规则奖励 GRPO 训练（Qwen3.5-27B ）}
\begin{itemize}
    \item 基于 verl 框架（Megatron + vLLM ），使用 6 万条病历质控数据对 Qwen3.5-27B 进行 GRPO 训练；
    \item 设计规则奖励函数，产出 0/1 奖励及 \texttt{judge\_reason} 分类标签（ok / mismatch / json\_not\_found 等）实时监控训练质量。
    \item 在 2{,}000 条验收集上 F1 值从基线 \textbf{87.4\%} 提升至 \textbf{89.27\%}；观察到精确率上升、召回率下降，分析为模型趋向保守输出；
    \item 进行不同训练数据难度配比实验和DAPO、Maxrl对照实验分析实验结果。
\end{itemize}

% -------- 病历生成 GenRM --------
\textbf{病历生成 GenRM GRPO 训练}
\begin{itemize}
    \item 从医患对话自动生成结构化病历（开放式任务，无标准答案）；Qwen3.5-27B 与 SFT 版本 LLM-as-Judge 评分均不理想，
    \item 设计生成式奖励模型评分维度，构建对应评测 prompt 与 rubric，以对齐人工评估标准；
\end{itemize}

\textbf{多任务联合训练}
\begin{itemize}
    \item 设计多任务联合训练方案，结合病历质控、KIE、医考与病历生成任务，提升模型在生成式任务上的表现；
    \item 通过多任务训练，模型在病历生成任务上 F1 值从基线 \textbf{82.3\%} 提升至 \textbf{85.6\%}。
% -------- 多模态评测 --------
\textbf{医疗多模态大语言模型自动评测与错误分析}
\begin{itemize}
    \item 设计涵盖多模态理解质量、画像结合度、医学专业性与准确性、逻辑推理、医疗安全与合规、回答实用性与可执行度等维度的自动评测框架；基于 LLM-as-a-Judge（Claude）结构化打分，完成 bad case 回溯与多维度评测报告可视化。
\end{itemize}

% -------- 数据构建 --------
\textbf{医疗垂域 SFT 数据构建与质检}
\begin{itemize}
    \item 整合开源病历和医院真实门诊数据，设计规则预筛+大模型双重质检双重流水线保障训练数据质量，基于deepseek-R1蒸馏Q-Cot-A,蒸馏强模型Q-A ,合成高质量慢思考和快思考SFT数据集。
\end{itemize}

% ====================== 个人总结 ======================
\section{个人总结}

\textbf{充足的实习时间:}可实习6个月\\[4pt]
\textbf{技术技能}
\begin{itemize}
    \item \textbf{大模型训练框架:}verl、Swift、Megatron-LM、vLLM、sglang
    \item \textbf{大模型评测与分析:}LLM-as-a-Judge、奖励函数设计、Rubric 设计、Bad Case 分析、Post-training Evaluation
    \item \textbf{工程与实验:}Python、PyTorch、Linux
\end{itemize}

\end{document}

% -------- 项目三（注释掉）：医考选择题 GenRM --------
% \textbf{医考选择题生成式奖励模型 GRPO 训练（进行中）}
% \begin{itemize}
%     \item \textbf{数据构造:}整理构造 12{,}950 条中文多选题与 65{,}293 条中英文单选题,将单选题筛选并改写为开放式问答,
%           最终保留 53{,}668 条高质量医考开放式 QA;组织多模型采样回答并用 LLM 自动打分,
%           构建训练集 20{,}000 query × 8 answers = 160{,}000 条偏好数据候选集与 4{,}000 条验证集。
%     \item \textbf{生成式奖励模型(GenRM)设计:}采用 Qwen2.5-7B-Instruct 作为 GenRM 后端。
%     \item \textbf{对比实验设计(计划):}规则匹配 baseline / 判别式 RM / 生成式 RM 三组对比。
% \end{itemize}

% -------- 项目四：环境搭建 --------